% Insert these two subsections within the existing appendix section
% Routing and fusion: protocols and supplementary results.
% Remove their old versions from experimental_diagnostics.tex to avoid
% duplicate labels. Keep the Circular internal-routing subsection there.
% Uses the existing bm, amsmath, booktabs, and graphicx packages and PDFs.

\subsection{Output-space fusion: conditional algebraic identities}
\label{app:fusion_consistency}

\paragraph{Assumptions and linear constraints.}
Let $\bm u_\alpha=\alpha\bm u_1+\beta\bm u_2$ and
$p_\alpha=\alpha p_1+\beta p_2$, where $\beta=1-\alpha$,
$\alpha\in[0,1]$ is constant in space and time within each prediction
window, and both candidates share a mesh, parameter, time grid, and
pressure gauge. For a common linear operator $D_h$,
$D_h\bm u_\alpha=\alpha D_h\bm u_1+\beta D_h\bm u_2$.
Consequently, zero discrete divergence, a zero-mean pressure gauge,
and identical affine boundary constraints are preserved if both
candidates satisfy them. This conditional statement does not establish
that the learned candidate fields satisfy these constraints.

\paragraph{Nonlinear residuals.}
For a fixed bilinear discretization $N_h$, define
$\delta\bm u=\bm u_1-\bm u_2$. Then
\begin{equation}
N_h(\bm u_\alpha,\bm u_\alpha)
=\alpha N_h(\bm u_1,\bm u_1)+\beta N_h(\bm u_2,\bm u_2)
-\alpha\beta N_h(\delta\bm u,\delta\bm u).
\end{equation}
Using the sign convention
$R_u(\bm u,p)=\partial_t\bm u-\bm f_h-L_h\bm u
+N_h(\bm u,\bm u)+G_hp$, this yields
\begin{equation}
R_u(\bm u_\alpha,p_\alpha)
=\alpha R_u(\bm u_1,p_1)+\beta R_u(\bm u_2,p_2)
-\alpha\beta N_h(\delta\bm u,\delta\bm u).
\end{equation}
Similarly, for
$R_p(\bm u,p)=A_hp-[\bm s_0+S_1\bm u+B_p(\bm u,\bm u)]$
with fixed bilinear $B_p$,
\begin{equation}
R_p(\bm u_\alpha,p_\alpha)
=\alpha R_p(\bm u_1,p_1)+\beta R_p(\bm u_2,p_2)
+\alpha\beta B_p(\delta\bm u,\delta\bm u).
\end{equation}
The additional terms separate the fusion defect from the candidates'
existing residuals; they need not increase or decrease the residual
norm. These are algebraic identities under the stated assumptions,
not evaluations of the native CFD equations. State-dependent limiters
or other non-bilinear discretizations require their actual operators.

\paragraph{Kinetic energy.}
For $E(\bm u)=\tfrac12\|\bm u\|_M^2$ with a common positive
quadrature matrix $M$,
\begin{equation}
E(\bm u_\alpha)=\alpha E(\bm u_1)+\beta E(\bm u_2)
-\frac{\alpha\beta}{2}\|\bm u_1-\bm u_2\|_M^2.
\end{equation}
Thus candidate disagreement reduces the blend's energy relative to the
weighted candidate energies, potentially including phase cancellation.
It does not imply lower energy than both candidates or improved agreement
with the reference. The blend is not fed back into either specialist,
so this term is not a recurrent dissipative modification of their dynamics.

\subsection{Square: measured physical-field diagnostics}
\label{app:fusion_measured}

\paragraph{Evaluation set and comparison fields.}
We evaluate both candidates, their frozen T2-C blend, and a common
POD-reconstructed reference on the same 16 held-out windows per Square
interface used for fusion evaluation ($K=24$, output spacing 4).
The candidates are Steady/Hopf for S--H and Periodic/Hopf for H--P.
Each interface contains two parameter values with eight overlapping
windows each. No window is selected by prediction quality, and overlapping
windows are not treated as independent experimental replications.
These measurements concern reconstructed fields, not errors against
unprojected CFD fields or residuals of the original OpenFOAM discretization.

\paragraph{Spatial operators and temporal differences.}
On the common 9,400-cell grid, quadratic local least-squares fits define
$D_x,D_y$, and the Laplacian $L$. The initial stencil contains 24 nearest
cells with independent coordinate scaling; it is expanded only if the
polynomial matrix is rank-deficient or its condition number exceeds
$10^6$. A second version starts from 36 neighbors to test sensitivity;
actual stencils contain at most 72 cells. Both versions differentiate
a manufactured quadratic polynomial with maximum absolute error below
$1.4\times10^{-10}$. This check verifies polynomial reproduction, not
accuracy on all reconstructed flow fields. We compute
\begin{align}
\widehat R_u&=\delta_t\bm u+N_D(\bm u,\bm u)
+\nabla_Dp-\mathrm{Re}^{-1}L\bm u,\\
\widehat R_p&=Lp+D_xN_{D,x}+D_yN_{D,y},
\end{align}
where $N_D(\bm u,\bm u)=uD_x\bm u+vD_y\bm u$ and
$N_{D,x},N_{D,y}$ are its two components. Centered temporal differences
and summary statistics use output steps 2--23.

\paragraph{Metrics and averaging.}
Residuals are summarized by area-weighted RMS over the interior wake
$6\leq x\leq16$, $1\leq y\leq3$ (2,627 cells). For a scalar or vector
residual $r$, this is
\begin{equation}
\|r\|_{\mathrm{RMS},\Omega_w}
=\left(\frac{\sum_{i\in\Omega_w}w_i\|r_i\|_2^2}
{\sum_{i\in\Omega_w}w_i}\right)^{1/2}.
\end{equation}
Kinetic energy uses quadrature over the whole domain. We distinguish
the energy error from the candidate-relative fusion deficit:
\begin{align}
\varepsilon_E&=100\frac{|E(\bm u)-E_{\rm ref}|}{E_{\rm ref}},\\
d_E&=100\frac{\alpha E(\bm u_1)+\beta E(\bm u_2)
-E(\bm u_\alpha)}{E_{\rm ref}},
\end{align}
where $E_{\rm ref}=E(\bm u_{\rm ref})$. Both quantities are percentages,
but only $\varepsilon_E$ measures agreement with reference energy.
Statistics average the evaluated steps and windows within each parameter,
then weight the two parameters equally.

\begin{table}[tbp]
\centering\small
\caption{Square physical-field diagnostics with the 24-neighbor initial
stencil. Energy error is evaluated over the whole domain; residuals are
interior area-weighted RMS, not percentages. Entries average output steps
2--23 and the evaluation windows as described above. Reference denotes
POD-reconstructed CFD rather than raw CFD.}
\label{tab:fusion_residual_measured}
\begin{tabular}{@{}llrrr@{}}
\toprule
Interface & Field & Energy error (\%) & $\|\widehat R_u\|$ & $\|\widehat R_p\|$\\
\midrule
S--H & Reference & 0 & 0.00460 & 0.03395\\
& Steady & 2.81100 & 0.00530 & 0.03674\\
& Hopf & 0.14296 & 0.00607 & 0.03513\\
& T2-C & 0.17641 & 0.00520 & 0.03421\\
\midrule
H--P & Reference & 0 & 0.11486 & 0.03736\\
& Periodic & 0.09163 & 0.11627 & 0.03785\\
& Hopf & 0.17298 & 0.11189 & 0.04047\\
& T2-C & 0.05974 & 0.11408 & 0.03631\\
\bottomrule
\end{tabular}
\end{table}

\paragraph{Measured comparisons and interpretation.}
Table~\ref{tab:fusion_residual_measured} shows different outcomes for
different diagnostics. On S--H, T2-C has smaller momentum and pressure
residuals than either candidate, but Hopf has a smaller energy error.
On H--P, T2-C has the smallest pressure residual and energy error among
the predictions, whereas Hopf has the smallest momentum residual.
The pressure-residual ordering also holds with the 36-neighbor initial
stencil, although changing the stencil changes residual magnitudes
substantially. The reference residual is nonzero and includes POD
truncation and spatial/temporal differentiation error; it is not a solver
convergence tolerance. In particular, a predicted field can have a lower
diagnostic residual than the reference without being more accurate.

\paragraph{Algebraic verification versus physical accuracy.}
Across both stencils and all windows, the maximum absolute discrepancies
in the energy, momentum-residual, and pressure-residual identities are
$1.2\times10^{-15}$, $7.4\times10^{-15}$, and $6.1\times10^{-14}$,
respectively. These verify the implemented identities, not conservation
or physical accuracy. The mean normalized fusion deficit $d_E$ is
$0.00314\%$ for S--H and $0.00181\%$ for H--P; these values are not
the reference-energy errors in Table~\ref{tab:fusion_residual_measured}.
Figure~\ref{fig:fusion_energy_defects_supp} separates actual kinetic
energy from candidate-relative fusion defects.

\paragraph{Probe and phase diagnostics.}
A transverse-velocity probe is fixed at the cell nearest $(8,2.5)$.
We detrend each signal and compare Hilbert phases with the reference.
A window is marked unresolved if the reference dominant non-Nyquist
Fourier component spans fewer than two cycles or the detrended reference
standard deviation is below $10^{-5}$. No S--H window and 15 of the 16
H--P windows pass this screening. Passing it does not exclude temporal
aliasing or Hilbert endpoint effects. Consequently, these traces do not
establish a general phase-error improvement or long-time phase stability.
Figures~\ref{fig:square_sh_diagnostics_supp}
and~\ref{fig:square_hp_diagnostics_supp} show aggregate diagnostics and
the first frozen window's probe and phase traces. The displayed window
is unresolved in both figures; its phase trace is descriptive only,
not a resolved phase-accuracy comparison.

\begin{figure}[tbp]
\centering
\includegraphics[width=\linewidth]{figures/round3_sh_diagnostics.pdf}
\caption{Square S--H diagnostics. The first four panels average all 16
windows; the final two show the first frozen window rather than a
best-case selection. Candidate 1 is Steady and candidate 2 is Hopf;
the reference is POD-reconstructed CFD. The displayed phase trace fails
the stated resolution screen and is descriptive only.}
\label{fig:square_sh_diagnostics_supp}
\end{figure}

\begin{figure}[tbp]
\centering
\includegraphics[width=\linewidth]{figures/round3_hp_diagnostics.pdf}
\caption{Square H--P diagnostics under the same protocol. Candidate 1
is Periodic and candidate 2 is Hopf. The first four panels average all
16 windows; the probe and phase panels show the first frozen window.
Although 15 windows pass the phase screen, this displayed window does
not, so its phase trace is not evidence of resolved phase accuracy.}
\label{fig:square_hp_diagnostics_supp}
\end{figure}

\begin{figure}[tbp]
\centering
\includegraphics[width=.95\linewidth]{figures/round3_sh_energy_excess.pdf}
\includegraphics[width=.95\linewidth]{figures/round3_hp_energy_excess.pdf}
\caption{Kinetic energy and additional fusion defects for Square S--H
(upper block) and H--P (lower block), averaged over all windows.
The energy deficit is relative to the weighted candidate energies and
normalized by reference energy, not the error against reference energy.
Momentum and pressure excesses are differences from the weighted
candidate residuals, not the total fused residuals.}
\label{fig:fusion_energy_defects_supp}
\end{figure}
